% GSR: original global quotient, exact resolvent and full spectral receiver.
\section{The original global quotient, its resolvent, and the complete primary spectral receiver}
\label{sec:global-spectral-quotient-resolvent}

\paragraph{Human sources and precisely read mathematical inputs.}
The original adelic image, closed quotient, and trace are from
Alain Connes, Caterina Consani and Matilde Marcolli,
\href{https://arxiv.org/abs/math/0703392v1}{\emph{The Weil proof and
the geometry of the adeles class space}}, original author source
\nolinkurl{Yuri70v2.tex}, SHA-256
\nolinkurl{b1f694cb934dd4fef829287f1a55dd8de21c67a481e552453c2f7e8fe0ae6fb5}.
For this continuation the complete passages at lines 1499--1549,
1621--1681, and 1680--1831 were read, including
\texttt{motH1AKCK}, the following closure convention,
\texttt{specrealH1mot}, \texttt{rangecVdef},
\texttt{varthetaH1V}, \texttt{varthetagTrH1}, and \texttt{vanishV}.
The explicit formula and sharp conventions are retained from
GTK1--38, SWQ1--28, and WQP1--24.

Ralf Meyer,
\href{https://arxiv.org/abs/math/0311468v3}{\emph{On a representation
of the idele class group related to primes and zeros of
$L$-functions}}, original author source
\nolinkurl{original_author.tex}, SHA-256
\nolinkurl{9453236959ee2fc232386a203637a12ca8b7f801469d86a6af1ccd7f21800f05},
was read at lines 421--470, 1150--1244, 4150--4200,
4371--4631, and 4891--4926. In particular the exact locators are
\texttt{lem:Lap\_weighted\_Sch}, \texttt{the:difference\_spectrum},
\texttt{sec:global\_spectral\_analysis},
\texttt{pro:quasi\_nilpotent\_character}, and
\texttt{pro:split\_off\_quasi\_nilpotent}.
The original source retains additional vertical growth conditions
in its summation image and expressly does not settle whether its
zero representation has a quasi-nilpotent part. A trace identity
is therefore not used below as a proof that an entire quotient
equals a Hilbert direct sum.

The preceding programme finite-jet proof WJI1--18 was read in full.
The present construction differs in an essential respect: its
cardinal functions vanish at \emph{all} other original zeros,
and its resolvent acts on the actual closed global quotient.
All new maps and their topology are proved below. The original
source nuclear trace theorem remains the stated human input.

\subsection{Original radial algebra and an exact description of its adelic image}

Let $B_0$ be the original radial strong weighted smooth algebra
of GTK14, with
\[
 p_{A,m,j}(f)=\sup_{u\in\mathbb R}e^{A|u|}(1+|u|)^m
                         |\partial_u^j f(e^u)|.
\]
Its measure and operations remain
$dy/y$, multiplicative convolution, and
$f^\sharp(y)=y^{-1}\overline{f(y^{-1})}$.
Write $F(s)=Mf(s)=\int_0^\infty f(y)y^sdy/y$.
The transform is entire, and on each bounded vertical strip
it and every derivative decrease faster than every power of
$|\Im s|$, by integration by parts in $u$.
Conversely this uniform vertical-strip Schwartz condition
gives an element of $B_0$ by Fourier inversion and contour
translation. This is Meyer's cited
\texttt{lem:Lap\_weighted\_Sch}; here the forward estimates
and every division used below are also proved directly in
the original $y$ variable.

Let $V_0$ be the compact-unit-invariant part of the actual
image $V=E\mathcal S(\mathbb A_{\mathbb Q})_0$.
Its exact radial description is
\begin{equation}\tag{GSR1}
 V_0=\left\{\,y\longmapsto2\sum_{n\ge1}\phi(ny):
       \phi\in\mathcal S(\mathbb R),\ \phi(-x)=\phi(x),\
       \phi(0)=0,\ \int_{\mathbb R}\phi(x)dx=0\,\right\}.
\end{equation}
Here is a proof of both directions. Average the original
adelic antecedent over the compact unit group of Haar mass
one. This preserves its value at zero and its additive
integral, since every unit dilation has additive modulus one.
A finite-place Schwartz function invariant under units is a
finite linear combination of characteristic balls
$1_{p^k\mathbb Z_p}$: its valuation profile is zero below
a lower bound, constant above an upper bound by local
constancy at zero, and has only finitely many intervening
values. Differences of successive ball indicators give
that profile exactly. Tensor expansion at the finitely
many exceptional primes consequently expresses the
antecedent as a finite sum
$\phi_j(x_\infty)\prod_p1_{p^{k_{j,p}}\mathbb Z_p}(x_p)$.
Replace $\phi_j$ by its even part, which leaves the
full rational-sign reduction unchanged.

Put $N_j=\prod_pp^{k_{j,p}}>0$. Its reduction is
$2\sum_{n\ge1}\phi_j(N_jny)$. Thus
$\phi(x)=\sum_j\phi_j(N_jx)$ gives exactly GSR1.
Its value at zero is the original adelic value at zero,
and its integral is
$\sum_jN_j^{-1}\int\phi_j$, exactly the original finite
Haar factors. Both are zero. Conversely any $\phi$
in GSR1 gives the actual antecedent
$\phi(x_\infty)\prod_p1_{\mathbb Z_p}(x_p)$ with the
two required zero moments, proving the converse.

Set
\begin{equation}\tag{GSR2}
 W_0=\overline{V_0}^{\,B_0},\qquad
 \mathcal C=B_0/W_0,\qquad q:B_0\longrightarrow\mathcal C,\qquad
 A f=-y\partial_y f .
\end{equation}
The continuous compact-unit projection proves that
$W_0=B_0\cap\overline V$ in the original whole Bruhat
space: apply that projection to an approximating net.
Thus $\mathcal C$ is the actual original-zeta sector of
the closed quotient, not a new closure convention.
It is a Hausdorff complete Fréchet space. By GSR1,
$A$ preserves $V_0$: on the even antecedent it is
$-x\phi'(x)$, whose value at zero is zero and whose
integral is $\int\phi=0$ by integration by parts.
Termwise differentiation commutes with $E$ by Schwartz
bounds. Continuity of $A$ therefore proves
$A(W_0)\subset W_0$, giving the continuous operator
$\bar A$ on $\mathcal C$.

Retain exactly the original theta-image test of GTK2,
\begin{equation}\tag{GSR3}
 \begin{gathered}
 \eta(x)=\pi x^2(\pi x^2-\tfrac32)e^{-\pi x^2},\qquad
 h(y)=2\sum_{n\ge1}\eta(ny)\in V_0,\\
 M(s)=Mh(s)=2\zeta(s)\left[
 \frac{\pi^2}{2\pi^{(s+4)/2}}\Gamma((s+4)/2)
 -\frac{3\pi}{4\pi^{(s+2)/2}}\Gamma((s+2)/2)\right],\\
 M(0)=M(1)=\tfrac14,\qquad M(s)=M(1-s).
 \end{gathered}
\end{equation}
Its entire interpretation, endpoint contributions, and the
retained multiplier comparison are GTK6--9. The original
zeta functional equation with that full multiplier, the
Euler product on $\Re s>1$, and the trivial-zero/pole
germs in ZTC17--17a show that the zeros of $M$ are
precisely the original nontrivial zeta zeros, with the
same multiplicities. Denote the distinct zeros by
$\mathcal Z$, their original multiplicities by $m_\rho$,
and put $\sigma(\rho)=1-\bar\rho$. Then
$m_{\sigma(\rho)}=m_\rho$. The zero count
$\sum_{|\rho|\le R}m_\rho=O(R\log(R+2))$ is proved
directly from this $h$ by Jensen's formula in GTK29.
Infinitude of the original nontrivial zero set is the
classical zeta fact retained by the source spectral
realization; no finite-height list replaces it here.

\subsection{Division in the original test space and global cardinal primary tests}

For $g\in B_0$ with $Mg(z)=0$, define
\begin{equation}\tag{GSR4}
 \mathcal D_zg(y)=
    -y^{-z}\int_0^y x^{z-1}g(x)\,dx
     =y^{-z}\int_y^\infty x^{z-1}g(x)\,dx .
\end{equation}
The equality is the actual vanishing Mellin moment,
not the omission of an endpoint. Every integral is
absolute. In $u=\log y$ coordinates, use the first
integral at $u<0$ and the second at $u>0$. For any
fixed compact set of $z$ and any requested exponential
weight, choose a larger weight in $g$; integration of
the resulting decaying exponential bounds both tails.
Differentiating the formula gives
$(A-z)\mathcal D_zg=g$, so its higher derivatives
are bounded recursively by those of $g$ and the
zeroth-order estimate. Parameter derivatives insert
powers of $u-v$ in the integrals, absorbed by a
slightly larger exponential weight. These estimates
prove $\mathcal D_zg\in B_0$, continuity on the
zero-moment subspace, and holomorphic dependence
for a holomorphic family $g_z$ with $Mg_z(z)=0$.
Integration by parts, with the now-proved vanishing
boundary terms, gives
\begin{equation}\tag{GSR5}
 M(\mathcal D_zg)(s)=\frac{Mg(s)}{s-z},
\end{equation}
including its removable value at $s=z$.

At $\rho\in\mathcal Z$, successive divisions are
permitted exactly $m_\rho$ times for $h$:
the first quotient has a zero of order $m_\rho-1$
there, and so on. Define
$h_\rho=\mathcal D_\rho^{m_\rho}h$ and
$u_\rho(t)=M(\rho+t)/t^{m_\rho}$, with
$u_\rho(0)\ne0$. For $0\le k<m_\rho$ let
$P_{\rho,k}(t)$ be the Taylor polynomial of
$t^k/u_\rho(t)$ modulo $t^{m_\rho}$, with its
ordinary Taylor coefficients. The actual tests
\begin{equation}\tag{GSR6}
 b_{\rho,k}=P_{\rho,k}(A-\rho)h_\rho\in B_0
\end{equation}
then have the exact full-primary jets
\begin{equation}\tag{GSR7}
 \left[Mb_{\rho,k}(\nu+t)\right]_{t^{m_\nu}}
  =\begin{cases}t^k,&\nu=\rho,\\0,&\nu\ne\rho.\end{cases}
\end{equation}
Indeed their transform is
$M(s)(s-\rho)^{-m_\rho}P_{\rho,k}(s-\rho)$.
At $\rho$ its jet is the prescribed inverse Taylor
product, while at every other zero $\nu$ it retains
the factor of order $m_\nu$. Multiplication by
the finite polynomial is exactly the original
differential operator in GSR6. Thus this is
global cardinal interpolation, annihilating all
other zeros with all their multiplicities.
In particular $Mb_{\rho,0}(\nu)=1_{\nu=\rho}$.

\subsection{Full jets, the actual trace radical, and the image topology}

Define
\begin{equation}\tag{GSR8}
 \begin{gathered}
 \mathcal Jf=
  \left(\left[Mf(\rho+t)\right]_{t^{m_\rho}}\right)_\rho,
 \qquad
 \mathcal Ef=(Mf(\rho))_\rho,\\
 N=\ker\mathcal J,\qquad R=\ker\mathcal E,\qquad
 \mathcal I_{\rm jet}=\mathcal J(B_0),\qquad
 \mathcal I_{\rm val}=\mathcal E(B_0).
 \end{gathered}
\end{equation}
Equip the two images with the exact quotient
topologies $B_0/N$ and $B_0/R$. Each individual
Taylor coefficient is a continuous functional
on $B_0$, so both kernels are closed and these
are Hausdorff Fréchet image spaces. This
specifies their topology; it does not silently
replace them by product or Hilbert subspace
topologies.

The actual source inclusion is
\begin{equation}\tag{GSR9}
 W_0\subset N\subset R.
\end{equation}
For the first inclusion, unfold GSR1 on $\Re s>1$:
$M(E\phi)(s)=2\zeta(s)\int_0^\infty\phi(x)x^{s-1}dx$.
The second factor is holomorphic for $\Re s>0$
because $\phi$ is bounded at zero and Schwartz
at infinity. It therefore gives the original
zeta zero of full order at each $\rho$.
Continuity of every jet passes this to the
closure $W_0$. The second inclusion is
the zeroth Taylor coefficient.

The complete source trace pairing is exactly
\begin{equation}\tag{GSR10}
 Q(f,g)=\sum_{\rho\in\mathcal Z}
       m_\rho Mf(\rho)\overline{Mg(\sigma(\rho))}.
\end{equation}
This is the author's trace theorem applied to
the original sharp, with all its global
endpoint, Gamma, real and prime terms retained
by GTK13 and GTK21. The sum is absolutely
convergent by vertical-strip decay and GTK29.
Its radical on $B_0$ is \emph{exactly} $R$.
One inclusion follows immediately from the
formula. Conversely, pairing $f$ with the
actual cardinal test $b_{\sigma(\rho),0}$
gives $m_\rho Mf(\rho)$. If $f$ pairs to zero
with every test, each of these values is zero.
This proves the reverse inclusion without a
spectral completeness assumption.

Consequently the radical on the actual
closed quotient $\mathcal C$ is $R/W_0$,
and the full sequence of proved receiving
maps and their kernels is
\begin{equation}\tag{GSR11}
 \mathcal C=B_0/W_0
    \longrightarrow\mathcal I_{\rm jet}=B_0/N
    \longrightarrow\mathcal I_{\rm val}=B_0/R,
 \quad
 \ker_1=N/W_0,\quad \ker_2=R/N .
\end{equation}
Both maps are onto their specified images.
By GSR7 the jet image contains the entire
algebraic direct sum
$\bigoplus_\rho\mathbb C[t]/(t^{m_\rho})$.
In particular $R/N$ contains every finite
collection of nilpotent primary directions
$\bigoplus_\rho t\mathbb C[t]/(t^{m_\rho})$,
with their exact coefficients.
The trace kernel cannot identify those
directions with the original closed-image
kernel. This statement includes the case
$m_\rho=1$, when that particular nilpotent
summand is zero.

\subsection{A resolvent on the actual closed quotient, not only on numerical values}

For $z\notin\mathcal Z$ set
\begin{equation}\tag{GSR12}
 \begin{gathered}
 R_zf(y)=y^{-z}\int_0^y x^{z-1}
       \left[f(x)-\frac{Mf(z)}{M(z)}h(x)\right]dx,\\
 M(R_zf)(s)=
      \frac{Mf(s)-Mf(z)M(s)/M(z)}{z-s}.
 \end{gathered}
\end{equation}
The numerator has zero moment at $z$, so
GSR4--5 prove $R_zf\in B_0$ and continuity.
They also prove its holomorphic dependence
on $z$ off $\mathcal Z$, and its meromorphic
dependence on the entire plane with possible
pole order at most $m_\rho$ at $\rho$.
More explicitly divide the holomorphic
zero-moment family $M(z)f-Mf(z)h$ first,
then divide its $B_0$-valued result by $M(z)$.
This proves the asserted meromorphy in
the actual strong topology.

It is essential to prove $R_z(W_0)\subset W_0$;
formal values at zeros do not prove that.
First take $\Re z>1$ and $f=E\phi\in V_0$
in the exact form GSR1. The corrected
antecedent
$\psi=\phi-Mf(z)M(z)^{-1}\eta$ is even
Schwartz, has value and additive integral
zero, and satisfies
$\int_0^\infty\psi(x)x^{z-1}dx=0$.
The last assertion follows by unfolding
both transforms and using the nonzero
Euler product $\zeta(z)$ in $\Re z>1$.
Put
\begin{equation}\tag{GSR13}
 \widetilde\psi(x)=x^{-z}\int_0^x
                      t^{z-1}\psi(t)dt\quad(x>0)
\end{equation}
and extend it evenly to $\mathbb R$.
At zero, the even Taylor expansion of
$\psi$ gives coefficients divided by
$z+2j$, none zero for $\Re z>1$.
Taylor's remainder under the integral
proves smoothness to every order, and
$\widetilde\psi(0)=\psi(0)/z=0$.
At infinity the zero moment rewrites
the integral as the negative tail,
giving rapid decrease with all derivatives.
Thus $\widetilde\psi$ is even Schwartz.
The identity
$(z+x\partial_x)\widetilde\psi=\psi$,
integrated with its actual vanishing
boundary terms, gives
$(z-1)\int_{\mathbb R}\widetilde\psi=\int\psi=0$.
It has both required source moments.
Hence $E\widetilde\psi\in V_0$.
Its transform is the second line of
GSR12, so Mellin injectivity gives
$R_zf=E\widetilde\psi$.
This proves $R_z(V_0)\subset V_0$
through an actual antecedent.

Continuity passes this result to $W_0$.
For fixed $w\in W_0$ the map
$z\mapsto q(R_zw)$ is holomorphic on
the connected domain $\mathbb C\setminus
\mathcal Z$ and zero on $\Re z>1$.
Applying each continuous linear
functional on the Hausdorff locally
convex quotient, the identity theorem
and separation of points give zero
throughout that domain. Thus the
preservation of $W_0$ holds for
every $z\notin\mathcal Z$, proving
the descended actual operator
$\bar R_z$.
The original differential equations
give the two inverse identities
\begin{equation}\tag{GSR14}
 \begin{gathered}
 (z-A)R_zf=f-\frac{Mf(z)}{M(z)}h,\qquad
 R_z(z-A)f=f,\\
 (z-\bar A)\bar R_z=\bar R_z(z-\bar A)=1_{\mathcal C}.
 \end{gathered}
\end{equation}
For the second equation the corrected
moment of $(z-A)f$ at $z$ is zero;
the kernel of $z-A$ in $B_0$ is
zero, since its solutions are
constant multiples of $y^{-z}$,
which cannot have both required
end decays. This proves that the
resolvent belongs to the actual
quotient, not an assumed spectral
replacement.

On the full jet image the exact
intertwining is
\begin{equation}\tag{GSR15}
 \mathcal J(R_zf)_\rho
    =(z-\rho-t)^{-1}\mathcal Jf_\rho
    =\sum_{k=0}^{m_\rho-1}
       \frac{t^k}{(z-\rho)^{k+1}}\mathcal Jf_\rho .
\end{equation}
The subtracted original $M(s)$ in
GSR12 vanishes to full order at
every zero, which proves this formula.
At a fixed $\rho$, the negative
Laurent coefficients of $R_zf$
have range spanned by the first
$m_\rho$ parameter derivatives of
$y^{-z}\int_0^y x^{z-1}h(x)dx$
at $\rho$. These functions belong
to $B_0$: their potentially
nondecaying tail coefficients are
the first $m_\rho$ zero moments
of $h$, all zero. Thus the
entire principal part has rank
at most $m_\rho$.

Its contour projection
$\Pi_\rho=(2\pi i)^{-1}\int\bar R_z\,dz$
on a small positive circle about
$\rho$ has rank exactly $m_\rho$.
Indeed GSR15 and GSR7 show that
its jet map covers that whole
primary block, giving rank at
least $m_\rho$; the preceding
range estimate gives the reverse
inequality. The resolvent identity,
obtained by multiplying the two
inverse equations in GSR14, and
Cauchy integration prove
$\Pi_\rho^2=\Pi_\rho$ and
$\Pi_\rho\Pi_\nu=0$ for $\rho\ne\nu$.
On its range $\mathcal J$ is an
isomorphism to
$\mathbb C[t]/(t^{m_\rho})$ and
$\bar A$ is exactly $\rho+t$.
Thus its actual pole order is
$m_\rho$, and its principal part is
\begin{equation}\tag{GSR16}
 \sum_{k=0}^{m_\rho-1}
     \frac{(\bar A-\rho)^k\Pi_\rho}
                             {(z-\rho)^{k+1}}.
\end{equation}
In particular the spectrum of the
actual quotient generator is
exactly $\mathcal Z$: GSR14 excludes
every other point and GSR16 makes
each indicated point a genuine pole.

There is also an exact description
of the kernel left in GSR11.
For $f\in N$, the ratio $Mf(z)/M(z)$
is entire because all zero orders
have been retained. GSR12 therefore
extends as an entire $B_0$-valued
function, sends $N$ into $N$, and
satisfies GSR14 on $N/W_0$ for
every $z\in\mathbb C$. Hence
\begin{equation}\tag{GSR17}
 N/W_0\text{ carries the restriction of }\bar A
       \text{ with an entire resolvent.}
\end{equation}
This constructs the exact remaining
subquotient and its operators.
It does not identify it with zero.
Meyer's explicit quasi-nilpotent
qualification and vertical growth
conditions are therefore preserved;
they are not overwritten by equality
of the trace characters.

\subsection{The entire direct-sum receiver and its exact indefinite form}

For the value pairing use
\begin{equation}\tag{GSR18}
 \mathscr H_{\rm val}
    =\{a:\sum_\rho m_\rho|a_\rho|^2<\infty\},
 \qquad
 [a,b]=\sum_\rho m_\rho a_\rho
                              \overline{b_{\sigma(\rho)}} .
\end{equation}
This is a continuous Hermitian form:
the permutation $Jb_\rho=b_{\sigma(\rho)}$
is a self-adjoint unitary for its
displayed positive Hilbert norm,
and $[a,b]=\langle a,Jb\rangle$.
No assertion that $[a,a]$ is
nonnegative is made. The original
test map $\mathcal E$ is continuous
into this space by vertical-strip
decay and the zero count. GSR7
puts every finitely supported
vector in its image, proving
that its Hilbert closure is all
of $\mathscr H_{\rm val}$.
This is the Hilbert completion of
the value receiver in this
explicit norm, not a claim
that $\mathcal C$ itself carries
that faithful norm.

The complete decomposition under
$\sigma$ has the original blocks
\begin{equation}\tag{GSR19}
 \begin{cases}
 m_\rho a_\rho\bar b_\rho,
        &\rho=\sigma(\rho),\\
 m_\rho(a_\rho\bar b_{\sigma(\rho)}
       +a_{\sigma(\rho)}\bar b_\rho),
        &\{\rho,\sigma(\rho)\}\text{ has two elements}.
 \end{cases}
\end{equation}
The latter Gram matrix is
$m_\rho\begin{pmatrix}0&1\\1&0\end{pmatrix}$
in the original two coordinates.
For example the actual global
test $b_{\rho,0}-b_{\sigma(\rho),0}$
has value $-2m_\rho$ in that
two-element case and no other
spectral contribution, by GSR7.
This proves the exact relation of
each permitted off-line orbit to
the original full form; it neither
asserts such an orbit exists nor
replaces it by sampled zeros.

The full primary receiver instead uses
\begin{equation}\tag{GSR20}
 \mathscr H_{\rm jet}=
   \mathop{\bigoplus\nolimits_{\rho\in\mathcal Z}}^{\ell^2}
       \mathbb C[t]/(t^{m_\rho}),
 \quad
 \|(a_{\rho,k})\|^2=\sum_\rho\sum_{k=0}^{m_\rho-1}|a_{\rho,k}|^2,
 \quad D_\rho=\rho I+N_\rho,\quad N_\rho a=t a .
\end{equation}
The Taylor coefficients here are
$Mf^{(k)}(\rho)/k!$; no factorial
has been suppressed.
Cauchy's formula on a circle of
radius one about $\rho$ bounds
every such coefficient by the
same supremum of $Mf$ on that
circle. Uniform vertical-strip
decay on the enlarged strip
therefore shows, for every $r$,
\[
 \sum_{\rho,k}(1+|\rho|)^{2r}
            |Mf^{(k)}(\rho)/k!|^2<\infty .
\]
Its estimate is continuous in
finitely many seminorms of $f$.
GSR7 supplies every finitely
supported jet vector; hence the
Hilbert closure of
$\mathcal I_{\rm jet}$ is exactly
$\mathscr H_{\rm jet}$. Its image
has the additional proved rapid
weighted decay and its original
quotient topology. Neither this
image nor that topology is
silently replaced by its closure.

Multiplication of full jets and
their original sharp are exactly
\begin{equation}\tag{GSR21}
 (ab)_\rho(t)=a_\rho(t)b_\rho(t)\pmod{t^{m_\rho}},
 \qquad
 (a^\sharp)_{\rho,k}
              =(-1)^k\overline{a_{\sigma(\rho),k}}.
\end{equation}
They follow from the original
convolution Mellin product and
$M(f^\sharp)(s)=\overline{Mf(1-\bar s)}$.
The ordinary finite-block trace
of multiplication by
$a_\rho b_\rho^\sharp$ is
$m_\rho a_{\rho,0}
\overline{b_{\sigma(\rho),0}}$.
Its nilpotent directions have
zero trace. Thus GSR18 is the
actual trace form of the full
primary receiver, with precisely
its proved nilpotent radical.
The factor $m_\rho$ is a
dimension in this jet trace.
It is not the ordinary Hilbert
operator trace of multiplication
on the single weighted line in
$\mathscr H_{\rm val}$.

\subsection{Global trace-class smoothing and resolvent trace duality}

On $\mathscr H_{\rm jet}$ define
$D$ blockwise as in GSR20 on
$\{a:\sum_\rho\|D_\rho a_\rho\|^2<\infty\}$.
Since $\|N_\rho\|\le1$ in the
displayed coefficient norm,
this is a closed densely defined
operator. At $z\notin\mathcal Z$,
its resolvent is the entire
block formula of GSR15.
On the tail where $|z-\rho|\ge2$,
the geometric operator sum has
norm at most $2/|z-\rho|$.
Each remaining block is finite
dimensional. Thus the resolvent
is bounded, compact, and
Hilbert--Schmidt, because
\[
 \sum_\rho m_\rho(1+|\rho|)^{-2}<\infty
\]
by the retained zero count.
Truncating the direct sum proves
compactness directly.

For every original $f\in B_0$,
its integrated receiver is
\begin{equation}\tag{GSR22}
 \Gamma_f= Mf(D),\qquad
 (\Gamma_f)_\rho=
 \sum_{k=0}^{m_\rho-1}\frac{Mf^{(k)}(\rho)}{k!}N_\rho^k,
 \qquad
 \operatorname{Tr}\Gamma_f=\sum_\rho m_\rho Mf(\rho).
\end{equation}
This is trace class. Indeed the
block trace norm is at most
$m_\rho\sum_{k<m_\rho}|Mf^{(k)}(\rho)/k!|$.
The preceding radius-one Cauchy
bound is rapid in $|\rho|$,
while $m_\rho\le C(1+|\rho|)
\log(2+|\rho|)$ by the same
zero count. A sufficiently
large decay exponent makes
the sum of these bounds finite.
The trace formula then follows
by summing the diagonal of
each triangular block.
It is exactly the original
source trace of $f$, with its
full explicit formula as in
GTK13, not a newly inferred
positivity theorem.

The resolvent duality also
exists on the actual quotient
without assigning a trace to
its bare resolvent. Convolution
by $R_zf$ is
$\bar R_z$ followed by convolution
by $f$: multiply by
$z-\bar A$ and use GSR14 and
the fact that $h$ acts as zero.
Consequently the source
trace-class theorem gives
\begin{equation}\tag{GSR23}
 \begin{gathered}
 \operatorname{Tr}_{\mathcal C}
       \bigl(\bar R_z\,\underline\vartheta(f)\bigr)
  =\mathcal T(R_zf)
  =\sum_\rho\frac{m_\rho Mf(\rho)}{z-\rho}
  =\operatorname{Tr}\bigl((z-D)^{-1}\Gamma_f\bigr),\\
 Q(R_zf,g)=
   \sum_\rho\frac{m_\rho Mf(\rho)
                 \overline{Mg(\sigma(\rho))}}{z-\rho}.
 \end{gathered}
\end{equation}
Every series is absolute, locally
uniform off its poles, and
has a simple scalar pole of
residue $m_\rho Mf(\rho)$
or the displayed paired value.
The actual operator resolvent
still has order $m_\rho$ at
that point. Scalar trace and
operator pole orders therefore
record different original data.
The first equality uses the
original author's trace notion
on $\mathcal C$, and the last
uses an ordinary trace-class
Hilbert operator on the
explicit receiving space;
their equality is proved,
not an identification of the
underlying spaces.

On the value space the exact
indefinite adjoint is
$D^\dagger=JD^*J=1-D$.
Equivalently, with both variables
in their domains,
\begin{equation}\tag{GSR24}
 [Da,b]=[a,(1-D)b],\qquad
 \bigl((z-D)^{-1}\bigr)^\dagger=-(1-\bar z-D)^{-1}.
\end{equation}
In the first equality the
coefficient on the second
side is
$1-\overline{\sigma(\rho)}=\rho$.
The second follows by inversion
and retains its minus sign.
On primary jets the original
sharp in GSR21 gives
$D^\sharp=1-D$, including
$t\mapsto-t$. These are exact
weight-one duality equations
for the actual form; no
positive inner product has
been inferred from them.

For bare spectral resolvents
use the convergent difference,
not a divergent unregularized
sum. Products of the two
Hilbert--Schmidt resolvents are
trace class, and their identity
gives
\begin{equation}\tag{GSR25}
 \begin{gathered}
 \operatorname{Tr}\bigl((z-D)^{-1}-(z_0-D)^{-1}\bigr)
  =\sum_\rho m_\rho
       \left[\frac1{z-\rho}-\frac1{z_0-\rho}\right],\\
 \operatorname{Tr}(z-D)^{-2}
                   =\sum_\rho\frac{m_\rho}{(z-\rho)^2}.
 \end{gathered}
\end{equation}
The first summand is $O(|\rho|^{-2})$
and the second likewise; triangular
nilpotents have zero trace, but
remain in the full resolvents.
These traces are also limits of
actual source trace-class tests.
Retain
$f_\varepsilon(y)=(4\pi\varepsilon)^{-1/2}
\exp(-(\log y)^2/(4\varepsilon))$, $\varepsilon>0$.
The original Mellin integral
is $Mf_\varepsilon(s)=e^{\varepsilon s^2}$
by the Gaussian Fourier integral,
with its indicated Jacobian
$dy/y$. Apply GSR23 to the
resolvent difference and then
let $\varepsilon\downarrow0$.
For $0<\varepsilon\le1$ the
modulus of that factor at
$\rho$ is at most $e$ because
$0<\Re\rho<1$; dominated
convergence with the
$|\rho|^{-2}$ bound gives the
first trace in GSR25.
The same argument after
differentiation gives the
second. Thus the receiving
trace has an explicit limit
through original admissible
tests, without assigning an
unsupported trace to an
operator on $\mathcal C$.

\subsection{The exact infinite-rank obstruction for the diagonal tensor clock}

On the algebraic tensor product
of primary blocks retain
\begin{equation}\tag{GSR26}
 D^{[2]}_{\rho,\nu}
     =(\rho+\nu)I+N_\rho\otimes I+I\otimes N_\nu .
\end{equation}
Its block dimension is
$m_\rho m_\nu$. The nilpotent
has exact length
$m_\rho+m_\nu-1$:
its power $m_\rho+m_\nu-2$
on $1\otimes1$ is
$\binom{m_\rho+m_\nu-2}{m_\rho-1}
t^{m_\rho-1}\otimes t^{m_\nu-1}$,
which is nonzero, and the next
power is zero. Thus the
operator resolvent has this
full pole order, while its
block trace is only
$m_\rho m_\nu/(z-\rho-\nu)$.
For any original test $f$,
its integrated block is
the complete Taylor polynomial
of $Mf$ at $\rho+\nu$ in
that nilpotent, and its
block trace is
$m_\rho m_\nu Mf(\rho+\nu)$.
These statements are finite
algebra identities on every
block, with no claim that
their infinite traces converge.

They give an unconditional
global failure of that trace
class. The original functional
equation pairs every zero
$\rho$ with $1-\rho$, with
the same multiplicity.
On all infinitely many
ordered pairs $(\rho,1-\rho)$
the tensor eigenvalue is
exactly $1$. Retain the
original coefficient basis
vector
$v_{\rho,1-\rho}
=t^{m_\rho-1}\otimes t^{m_{1-\rho}-1}$.
It already has norm one
in GSR20, with no rescaling.
The tensor nilpotent kills
this vector, so it is an
eigenvector, and the vectors
from distinct blocks are
orthogonal in the specified
Hilbert tensor product.
For the actual original
theta-image test GSR3,
\begin{equation}\tag{GSR27}
 Mh(1)=\tfrac14,\qquad
 Mh(D^{[2]})v_{\rho,1-\rho}
                         =\tfrac14v_{\rho,1-\rho}.
\end{equation}
Hence its tensor-clock
operator cannot be compact
or trace class. The conclusion
already follows on this
orthogonal sequence; it
does not assume RH or
discard the nilpotents.
In fact its full blockwise
operator is bounded: on
the strip $0\le\Re s\le2$
use Cauchy's estimate for
$Mh$ on radius three;
the tensor nilpotent has
norm at most two, so the
Taylor-operator series is
bounded by a convergent
geometric series times
that strip supremum.
Thus noncompactness here
is a proved property of
a bounded full receiver,
not merely failure to
define a formal expression.

Likewise for any $z$ with
$\Re z<-3$, all block
resolvents are bounded
by a uniformly convergent
geometric series, but on
the displayed eigenvectors
they equal $(z-1)^{-1}$.
They are noncompact.
For tensor order $k\ge2$,
fix the remaining $k-2$
zeros: the same infinite
family has sum
$1+\rho_3+\cdots+\rho_k$.
This records the entire
infinite-rank obstruction
to carrying the one-variable
compact-resolvent argument
over to the diagonal tensor
clock.

These finite primary blocks
also embed in the actual
quotient by GSR16; the
independent original-coordinate
chain construction and its
continuous equivariant
projections are GDT13--17
in \nolinkurl{GLOBAL_DILATION_TENSOR_RECEIVER.tex}.
Their algebraic tensor
embeddings retain every
block above. The Hilbert
tensor completion used
for the stated compactness
calculation is the specified
receiving completion, and
is not identified with
the whole actual quotient
or an unspecified completion
of its tensor product.

Independent-variable
smoothing has a different
exact map:
$\Gamma_f\otimes\Gamma_g$
is trace class with trace
$\mathcal T(f)\mathcal T(g)$,
by trace-norm finite-rank
approximation and GSR22.
Its block eigenvalues are
$Mf(\rho)Mg(\nu)$, whereas
the diagonal-clock test
has $Mf(\rho+\nu)$.
Both formulas keep all
multiplicities and all
original Gamma and endpoint
contributions through GSR3.
In particular $\Gamma_h=0$
on every original primary
block, since $Mh$ vanishes
to order $m_\rho$ there;
but GSR27 is nonzero of
infinite rank. This is the
specific global trace and
quotient obstruction, not
an assertion that arbitrary
nonlinear products descend.

\subsection{The entire arbitrary-lattice receiver}

For the original nontrivial
bounded distributive lattice
$L$, all the linear maps
$q,\mathcal J,\mathcal E,A$,
$R_z$, $\Pi_\rho$ and their
displayed restrictions lift
by
\begin{equation}\tag{GSR28}
 (f,\lambda)\longmapsto(Tf,\lambda)
      \quad\hbox{on}\quad
 G_L(U)=\{(0,\lambda):\lambda\in L\}
                       \cup(U\times\{1_L\}).
\end{equation}
There is one label on each
entire original vector,
including its infinite
spectral image after that
amplitude has been proved
to converge. No infinite
lattice join or lattice
measure is introduced.
Linearity and distributivity
prove the semimodule maps
and their same-label
quotient fibres. In
particular the quotients
in GSR11 keep a zero
amplitude with every
original label, including
$e=(0,1_L)\ne\tau=(0,0_L)$.

The exact pairing receiver is
\begin{equation}\tag{GSR29}
 Q_L((f,\lambda),(g,\mu))
   =\left(\sum_\rho m_\rho Mf(\rho)
           \overline{Mg(\sigma(\rho))},
                                    \lambda\wedge\mu\right).
\end{equation}
Every numerical term has
its original global formula
GTK13 or GTK24. Zero
amplitude does not remove
the displayed meet.
On a primary block the
power $N_\rho^{m_\rho}$
lifts to $(v,\lambda)\mapsto
(0,\lambda)$, not the
constant-$\tau$ map.
The tensor receiver retains
one common label on its
whole amplitude tensor;
the infinite-rank sequence
in GSR27 lies at top
support, so its failure
of compactness is not
removed by any label
identification. At an
actual support prime $J_L$,
the existing localization
$(v,\lambda)\mapsto[\lambda]_{J_L}$
commutes with every displayed
label-preserving linear map
and takes the pairing to
the original meet. That
localization does not
assert an unchanged
arithmetic amplitude stalk.

The proved global relation is
therefore the actual quotient
and resolvent GSR12--17,
its full primary image and
specified closure GSR20,
its exact numerical radical
GSR10--11, and its original
trace-class receiving map
GSR22--25. The all-jet-vanishing
subquotient is retained with
its entire resolvent, as the
human sources require.
The diagonal tensor receiver
has the explicit infinite-rank
endpoint block GSR27.
No positive-kernel assumption,
identification of image with
closure, or finite-prime
certificate is used.
